\reviewsection{Communication, Behavior, and the Reliability of Adult Response}

\citet{peckhamhardin2015behavior} add another important dimension: actions that adults label challenging may be connected to a message that does not yet have a reliable pathway. Functional communication support begins by investigating the purpose, teaching an accessible alternative, and ensuring that partners honor the message. A break message that never produces a break teaches the learner that the new route is less reliable than the action adults are trying to replace. This builds on my submitted paper, \textit{From Compliance to Communication: Prevent--Teach--Reinforce Planning Through Access, Agency, and Richer Evidence}, without recreating its behavior plan \citep{garciabautista2026ptr}. That paper examined how allowing a familiar laptop to coexist with hands-on materials removed an unnecessary either-or condition. Module 3 adds the language and partner question: how could the learner say, ``keep this,'' ``different order,'' ``show me,'' ``not yet,'' or ``I meant something else'' before an adult completes the interpretation?

My posted \textit{AAC Review: TD Snap} made the implementation test concrete \citep{garciabautista2026tdsnap}. TD Snap offers multiple page sets, access routes, editing tools, bilingual Motor Plan options, and printable low-tech materials. Its value is not the number of features. Its value depends on feature matching, personalized vocabulary, backup access, partner practice, and whether the communicator can do more than request. In the fictional plan, Mateo would need Spanish--English vocabulary and family language; Eli would need robust multimodal access; and Luca would need easy access to health, privacy, fatigue, pacing, and repair messages. A classmate's review of the same tool added a useful EQUITAS question by placing customizable photographs beside the recurring subscription cost. Personal photographs can connect AAC to familiar people, activities, and interests; for a culturally and linguistically diverse student, personalization should also protect names, home-language vocabulary, family relationships, humor, and identity. Yet a feature is not accessible if a school can introduce it and the student or family cannot sustain it across home, summer, transitions, or a funding change. The team must decide who will pay, maintain an effective low-tech backup, and ask whether the student's voice remains available when the app or speech output is not.

\begin{figure}[!ht]
\centering
\newcommand{\commbox}[2]{%
  \fcolorbox{PrimaryRed}{SoftBlue}{%
    \parbox[c][0.68in][c]{1.45in}{\centering\textbf{#1}\\[-0.02in]\small #2}%
  }%
}
\commbox{Student Meaning}{comment, question, refusal, repair}
\hspace{0.01in}{\color{PrimaryRed}$\longrightarrow$}\hspace{0.01in}
\commbox{Access Route}{touch, gaze, gesture, symbol, text}
\hspace{0.01in}{\color{PrimaryRed}$\longrightarrow$}\hspace{0.01in}
\commbox{Partner Practice}{model, wait, notice, confirm, act}
\hspace{0.01in}{\color{PrimaryRed}$\longrightarrow$}\hspace{0.01in}
\commbox{Agency}{message changes what happens next}
\caption{My Module 3 access-to-agency synthesis. A device is one route inside a communication system; authorship depends on matched access and a partner response that lets the message matter.}
\label{fig:access-to-agency}
\end{figure}

\reviewsection{Puzzle Plan, EQUITAS, and Evidence Justice}

The Puzzle Plan is my interdisciplinary framework for bringing lived experience, education, data, research, technology, and advocacy together so that no diagnosis, observation, behavior, or score stands alone. EQUITAS is the equity-centered heart of that framework. It asks whose knowledge is included, what context is missing, how power and bias shape interpretation, and whether a support increases dignity, agency, belonging, and access. Through the Puzzle Plan and EQUITAS, communication support becomes evidence justice. When accessible modes are missing, schools can mistake speech, speed, or compliance for competence. A fair record would document the student's message, available modes, environmental conditions, partner response, breakdowns, successful repairs, family knowledge, and response to support. Even a home--school log can either increase or restrict power. A list of warnings, sitting, following directions, and good/fair/poor ratings tells adults about adult expectations. A student-centered record asks what the learner communicated, what worked, what barrier should change, and what the student wants shared.
